% ═══════════════════════════════════════════════════════════════
% MT-02b-ML: Musterlösung Minitest tener + Alter
% Spanisch Klasse 7
% ═══════════════════════════════════════════════════════════════
% Dauer: 10 Minuten
% Punkte: 28
% Inhalt: Zahlen 11-20, tener-Konjugation, Alter angeben
% ═══════════════════════════════════════════════════════════════

\documentclass[11pt, a4paper]{article}

\usepackage[utf8]{inputenc}
\usepackage[T1]{fontenc}
\usepackage[ngerman]{babel}
\usepackage{helvet}
\renewcommand{\familydefault}{\sfdefault}

\usepackage[margin=2cm]{geometry}
\usepackage{enumitem}
\usepackage{tabularx}
\usepackage{booktabs}

\usepackage{xcolor}
\usepackage{tcolorbox}
\tcbuselibrary{skins}

\usepackage{fancyhdr}
\usepackage{lastpage}
\usepackage{needspace}

% FARBEN
\definecolor{spanisch-color}{HTML}{c41e3a}
\definecolor{grau-color}{HTML}{888888}
\definecolor{hellgrau-color}{HTML}{f5f5f5}
\definecolor{loesung-color}{HTML}{c62828}

\colorlet{spanisch}{spanisch-color}
\colorlet{grau}{grau-color}
\colorlet{hellgrau}{hellgrau-color}
\colorlet{loesung}{loesung-color}

\newcommand{\fachfarbe}{spanisch}

\newcommand{\thema}{Minitest: tener + Alter -- MUSTERLÖSUNG}
\newcommand{\fach}{Spanisch}
\newcommand{\klasse}{7}
\newcommand{\maxpunkte}{28}
\newcommand{\zeit}{10 Minuten}
\newcommand{\datum}{\today}

\pagestyle{fancy}
\fancyhf{}
\renewcommand{\headrulewidth}{0.4pt}
\renewcommand{\footrulewidth}{0.4pt}

\lhead{\fach{} -- Klasse \klasse}
\chead{\textbf{MT-02b -- ML}}
\rhead{\datum}

\lfoot{\zeit}
\cfoot{}
\rfoot{Seite \thepage\ von \pageref{LastPage}}

\newcommand{\punktebox}[1]{\hfill\fbox{\small \_/#1}}
\newcommand{\luecke}[1]{\rule{#1}{0.4pt}}
\newcommand{\lsg}[1]{\textcolor{loesung}{\textbf{#1}}}

\newtcolorbox{infobox}{
    colback=hellgrau,
    colframe=\fachfarbe,
    boxrule=1pt,
    arc=0pt,
    left=8pt, right=8pt, top=8pt, bottom=8pt
}

\newenvironment{aufgabe}[1]{%
    \needspace{5\baselineskip}%
    \nopagebreak[4]%
    \vspace{10pt}
    \noindent\textbf{\textcolor{\fachfarbe}{Aufgabe #1}}
    \nopagebreak[4]%
    \vspace{5pt}
    \nopagebreak[4]%
    \begin{list}{}{\leftmargin=0pt}
    \item[]
}{%
    \end{list}
}

\newcommand{\es}[1]{\textcolor{\fachfarbe}{\textbf{#1}}}

\begin{document}

% ═══════════════════════════════════════════════════════════════
% KOPFBEREICH
% ═══════════════════════════════════════════════════════════════

\begin{center}
    {\Large\bfseries\textcolor{\fachfarbe}{\thema}}
\end{center}

\vspace{5pt}

\noindent
\begin{tabularx}{\textwidth}{|X|X|X|}
\hline
\textbf{Name:} \lsg{LÖSUNG} &
\textbf{Klasse:} \klasse &
\textbf{Datum:} \luecke{2cm} \\
\hline
\end{tabularx}

\vspace{10pt}

\begin{infobox}
\textbf{Zeit:} \zeit{} | \textbf{Punkte:} \maxpunkte{} | \textbf{Hilfsmittel:} keine
\end{infobox}

\vspace{15pt}

% ═══════════════════════════════════════════════════════════════
% AUFGABE 1: Zahlen zuordnen (10 Punkte)
% ═══════════════════════════════════════════════════════════════

\begin{aufgabe}{1 -- Zahlen zuordnen}
Ordne die spanischen Zahlen der richtigen Ziffer zu. \punktebox{10}
\end{aufgabe}

\begin{center}
\begin{tabular}{rcl}
\es{once} & \lsg{11} & 11 \\[0.6em]
\es{doce} & \lsg{12} & 12 \\[0.6em]
\es{trece} & \lsg{13} & 13 \\[0.6em]
\es{catorce} & \lsg{14} & 14 \\[0.6em]
\es{quince} & \lsg{15} & 15 \\[0.6em]
\es{dieciséis} & \lsg{16} & 16 \\[0.6em]
\es{diecisiete} & \lsg{17} & 17 \\[0.6em]
\es{dieciocho} & \lsg{18} & 18 \\[0.6em]
\es{diecinueve} & \lsg{19} & 19 \\[0.6em]
\es{veinte} & \lsg{20} & 20 \\
\end{tabular}
\end{center}

\vspace{15pt}

% ═══════════════════════════════════════════════════════════════
% AUFGABE 2: tener-Konjugation (8 Punkte)
% ═══════════════════════════════════════════════════════════════

\begin{aufgabe}{2 -- tener konjugieren}
Ergänze die richtige Form von \es{tener} (haben). \punktebox{8}
\end{aufgabe}

\noindent
a) Yo \lsg{tengo} 12 años. \hfill (ich habe)

\vspace{0.8em}

\noindent
b) Tú \lsg{tienes} un hermano. \hfill (du hast)

\vspace{0.8em}

\noindent
c) María \lsg{tiene} 15 años. \hfill (María hat)

\vspace{0.8em}

\noindent
d) Nosotros \lsg{tenemos} una casa grande. \hfill (wir haben)

\vspace{15pt}

% ═══════════════════════════════════════════════════════════════
% AUFGABE 3: Alter angeben (10 Punkte)
% ═══════════════════════════════════════════════════════════════

\begin{aufgabe}{3 -- Alter angeben}
Schreibe die Zahl als spanisches Wort. \punktebox{10}
\end{aufgabe}

\noindent
a) Mi hermano tiene \lsg{trece} años. \hfill (13)

\vspace{0.8em}

\noindent
b) Tengo \lsg{catorce} años. \hfill (14)

\vspace{0.8em}

\noindent
c) Mi madre tiene \lsg{quince} años. \hfill (15)

\vspace{0.8em}

\noindent
d) Mi abuelo tiene \lsg{diecisiete} años. \hfill (17)

\vspace{0.8em}

\noindent
e) Mi prima tiene \lsg{diecinueve} años. \hfill (19)

% ═══════════════════════════════════════════════════════════════
% PUNKTEÜBERSICHT
% ═══════════════════════════════════════════════════════════════

\vspace{25pt}
\hrule
\vspace{10pt}

\noindent
\begin{tabularx}{\textwidth}{|X|c|c|c||c|}
\hline
\textbf{Aufgabe} & 1 & 2 & 3 & \textbf{Gesamt} \\
\hline
\textbf{max. Punkte} & 10 & 8 & 10 & \textbf{28} \\
\hline
\textbf{erreicht} & \lsg{10} & \lsg{8} & \lsg{10} & \lsg{28} \\[0.8em]
\hline
\end{tabularx}

\vspace{10pt}

\begin{flushright}
\textbf{Note:} \lsg{14 NP}
\end{flushright}

\end{document}
